\section*{Hoofdstuk \ref{ch:b3:planetary-interiors-surfaces} --- Planetaire binnenwerelden en oppervlakken}

\begin{solution}{exo:b3:planetary-interiors-surfaces:1}
$\tfrac52C\bar\rho = 2.5\times0.35\times5429 = 4750$; $x^2 = (4750 -
3300)/(5429 - 3300) = 0.681$, $x = 0.825$, een kern van $0.825\times2440 =
\qty{2010}{km}$; $\rho_c = 3300 + 2129/0.825^3 = \qty{7090}{kg.m^{-3}}$.
\end{solution}

\begin{solution}{exo:b3:planetary-interiors-surfaces:2}
$m = 3500\times\tfrac43\pi(25)^3 = \qty{2.29e8}{kg}$; $E = \tfrac12 m v^2 =
\tfrac12\times\num{2.29e8}\times(\num{2.0e4})^2 = \qty{4.6e16}{J}$, of
$\num{4.6e16}/\num{4.2e15} = \qty{11}{Mt}$. De regel geeft een krater van
$20\times\qty{50}{m} = \qty{1}{km}$.
\end{solution}

\begin{solution}{exo:b3:planetary-interiors-surfaces:3}
$N_1 = 120/\num{100000} = \qty{1.2e-3}{km^{-2}}$, relatieve fout $1/\sqrt{120}
= \qty{9}{\percent}$. Op het lineaire deel, $T = N_1/c = \num{1.2e-3}/\num{4.875e-4}
= \qty{2.46}{Ga}$ (de exponentiële term, $\num{8e-7}$, is verwaarloosbaar), en
$\sigma_T = T/\sqrt{120} = \qty{0.22}{Ga}$.
\end{solution}

\begin{solution}{exo:b3:planetary-interiors-surfaces:4}
Oppervlakte $4\pi(\num{1.8215e6})^2 = \qty{4.17e13}{m^2}$; vermogen $2.5\times
\num{4.17e13} = \qty{1.04e14}{W} = \qty{104}{TW}$, 2.2 keer het totale vermogen van de Aarde
afkomstig van een oppervlak twaalf keer kleiner, en $2.5/0.092 = 27$ keer zijn
flux.
\end{solution}

\begin{solution}{exo:b3:planetary-interiors-surfaces:5}
(a)~$\bar\rho = \num{6.417e23}/(\tfrac43\pi(\num{3.3895e6})^3) =
\qty{3934}{kg.m^{-3}}$; $g = GM/R^2 = \qty{3.73}{m.s^{-2}}$; $\rho gR =
\qty{4.97e10}{Pa}$. (b)~$19\mu/(2\rho gR) = 19\times\num{6.7e10}/\num{9.94e10}
= 12.8$, $k_2 = 1.5/13.8 = 0.109$. (c)~De gemeten waarde ligt
2.6 standaarddeviaties hoger: Mars geeft meer dan een vaste planeet met
de starheid van een mantel, zoals het zou doen met een vloeibare kern erin.
\end{solution}

\begin{solution}{exo:b3:planetary-interiors-surfaces:6}
(a)~Maan: $4\times9.8/1.62 = \qty{24}{km}$ in kristallijn gesteente,
\qty{12}{km} in sedimenten; Mars en Mercurius: $4\times9.8/3.7 =
\qty{10.6}{km}$ en \qty{5.3}{km}. De \qty{20}{km} van de Maan ligt tussen
zijn twee voorspellingen: de regel klopt binnen enkele tientallen procenten. (b)~$Y =
\rho gD/3 = 2600\times9.8\times4000/3 = \qty{34}{MPa}$, de sterkte van
gebarsten gesteente in plaats van die van een intact kristal.
\end{solution}

\begin{solution}{exo:b3:planetary-interiors-surfaces:7}
Mars: $\num{3.22e-14}(e^{6.93T} - 1) + \num{4.875e-4}\,T = \num{2.0e-3}$;
de lineaire schatting $4.1$~Ga is veel te oud, want de exponentiële term is
daar groot; bisectie geeft $T = \qty{3.34}{Ga}$ (exponentiële term
$\num{3.6e-4}$, lineaire term $\num{1.63e-3}$). De lunaire functie geeft
\qty{2.39}{Ga}. Elke functie is geijkt voor één hemellichaam: de snelheid van
inslagen en de grootte van de kraters die een gegeven projectiel maakt, verschillen
tussen de Maan en Mars.
\end{solution}

\begin{solution}{exo:b3:planetary-interiors-surfaces:8}
(a)~$3.551181/1.769138 = 2.0073$ en $7.154553/3.551181 = 2.0147$, dicht bij
$2$. (b)~$0.565247 - 0.844789 + 0.279542 = 0.000000$ per dag, tot op
$\num{1e-7}$: de drie gemiddelde bewegingen zijn vergrendeld, een resonantie die
de relatieve posities van de manen bij elke conjunctie. (c)~De trekkrachten herhalen zich op
dezelfde punten van de baan, zodat ze optellen in plaats van uit te middelen;
ze voeden de excentriciteit net zo snel aan als de getijden haar dempen.
\end{solution}

\begin{solution}{exo:b3:planetary-interiors-surfaces:9}
$n = 2\pi/(1.769138\times86400) = \qty{4.111e-5}{s^{-1}}$, $nR =
\qty{74.9}{m.s^{-1}}$, $(nR)^5/G = \qty{3.53e19}{W}$;
$\tfrac{21}{2}e^2(nR)^5/G = 10.5\times\num{1.6e-5}\times\num{3.53e19} =
\qty{5.9e15}{W}$ per eenheid van $k_2/Q$. Dan $k_2/Q =
\num{1.04e14}/\num{5.9e15} = 0.018$.
\end{solution}

\begin{solution}{exo:b3:planetary-interiors-surfaces:10}
(a)~$x^2 = (\tfrac52C\bar\rho - \rho_m)/(\bar\rho - \rho_m)$ is lineair in
$C$, met helling $\tfrac52\bar\rho/(\bar\rho - \rho_m)$. (b)~De helling is
$2.5\times3340/52 = 161$: $x^2$ verandert met $0.161$ bij $0.001$ van $C$.
Voor $C = 0.395$, $x^2 = 0.036 + 0.161 = 0.197$, $x = 0.44$, $\rho_c =
\qty{3880}{kg.m^{-3}}$; voor $C = 0.393$, $x^2 < 0$: helemaal geen kern.
(c)~Nabij $C = \tfrac25$ voelt het traagheidsmoment nauwelijks een kleine kern:
zijn derde significante cijfer beslist tussen een grote, lichte kern en geen.
Seismische golven, getijden en het magnetisch veld zijn nodig.
\end{solution}

\begin{solution}{exo:b3:planetary-interiors-surfaces:11}
(a)~$\sigma_T = N_1/(10\,N_1')$: bij \qty{1.0}{Ga}, $N_1 \approx N_1' \approx
\num{4.875e-4}$, $\sigma_T = \qty{0.10}{Ga}$; bij \qty{3.8}{Ga}, $N_1 =
\num{1.07e-2}$ en $N_1' = \num{6.15e-2}$, $\sigma_T = \qty{0.017}{Ga}$,
zes keer kleiner. (b)~Met $N_1 = cT$, $N_1' = c$ en $\sigma_T =
cT/(\sqrt n\,c) = T/\sqrt n$. (c)~De chronologiefunctie zelf is
onzeker: ze berust op enkele gedateerde locaties op de Maan en op de schaling
van de Maan naar Mars, fouten die geen telling kan wegnemen.
\end{solution}

\begin{solution}{exo:b3:planetary-interiors-surfaces:12}
(a)~De verhouding is $(n_ER_E/n_IR_I)^5(e_E/e_I)^2 = (0.4982\times0.8569)^5
\times(0.009/0.004)^2 = 0.0142\times5.06 = 0.072$: ongeveer
\qty{7.5e12}{W}, of $\num{7.5e12}/(4\pi(\num{1.5608e6})^2) =
\qty{0.24}{W.m^{-2}}$. (b)~Een oceaan ontkoppelt de ijsschil van de
binnenkant: de schil buigt veel meer dan een vaste maan zou doen, en het ijs
dissipeert anders dan het hete gesteente van Io, dus $k_2$ en $Q$ veranderen beide.
\end{solution}

\begin{solution}{pb:b3:planetary-interiors-surfaces:1}
\textbf{1.} $\bar\rho = \num{6.4169e23}/(\tfrac43\pi(\num{3.3895e6})^3) =
\qty{3934}{kg.m^{-3}}$.

\textbf{2.} $\tfrac52C\bar\rho = 3600$; $x^2 = (3600 - 3500)/(3934 - 3500) =
0.230$, $x = 0.479$, een kern van \qty{1620}{km}; $\rho_c = 3500 +
434/0.479^3 = \qty{7450}{kg.m^{-3}}$.

\textbf{3.} $x = 1830/3389.5 = 0.540$; $\rho_m = (3600 - 0.2915\times
3934)/(1 - 0.2915) = \qty{3460}{kg.m^{-3}}$, $\rho_c = 3462 +
472/0.1574 = \qty{6460}{kg.m^{-3}}$.

\textbf{4.} Het is een derde lager. De druk in de kern van Mars is veel
lager dan in die van de Aarde, dus het metaal is minder samengedrukt; en
een zo lichte kern wijst op veel lichtere elementen, gemengd met het ijzer.

\textbf{5.} Het \omterm{def:b3:planetary-interiors-surfaces:love-number}{Lovegetal} (2003) dateert ruim vóór InSight (2021): het
toonde al aan dat ten minste de buitenkern vloeibaar is en plaatste de
straal tussen \qty{1520}{km} en \qty{1840}{km}, een bereik dat de
seismologische waarde bevat.

\textbf{6.} $1 - (1 - 20/3389.5)^3 = \qty{1.8}{\percent}$ tot $1 - (1 -
37/3389.5)^3 = \qty{3.2}{\percent}$ van het volume, en minder van de massa:
lichter dan de mantel verlaagt het de effectieve manteldichtheid een beetje,
een effect dat kleiner is dan de onzekerheid van het model.

\textbf{7.} Vanaf \qty{1}{km}: 54; vanaf 1.41: 16; vanaf 2: 10; vanaf 2.83: 5;
vanaf 4, 5.66, 8 en 11.3: 2; vanaf 16: 1.

\textbf{8.} $N = 54/80000 = \qty{6.75e-4}{km^{-2}}$, $\sigma = \sqrt{54}/80000
= \qty{0.92e-4}{km^{-2}}$ (\qty{14}{\percent}).

\textbf{9.} Verwacht $54\times0.355 = 19$ kraters vanaf \qty{1.41}{km} en
$54\times0.147 = 7.9$ vanaf \qty{2}{km}; geteld 16 en 10, binnen één
standaarddeviatie ($\sqrt{19} = 4.4$, $\sqrt{7.9} = 2.8$): consistent.

\textbf{10.} Nee: $0.072/\num{6.75e-4} = 107$ keer lager dan dat. De telling
registreert tijd.

\textbf{11.} Klassen van gelijke breedte in $\log D$ bestrijken het bereik
gelijkmatig en houden vergelijkbare aantallen kraters in de klassen van een steile verdeling.

\textbf{12.} Het zijn secundaire kraters, gevormd door blokken die uit een
grote primaire krater zijn weggeslingerd: geen onafhankelijke inslagen uit de ruimte; ze zouden
kraters toevoegen die de \omterm{def:b3:planetary-interiors-surfaces:crater-dating}{kraterchronologiefunctie} niet telt.

\textbf{13.} $T = \num{6.75e-4}/\num{4.875e-4} = \qty{1.385}{Ga}$; de
exponentiële term, $\num{3.22e-14}(e^{9.6} - 1) = \num{4.7e-10}$, is een
miljoenste van $N_1$.

\textbf{14.} $\sigma_T = T/\sqrt{54} = \qty{0.19}{Ga}$: $1.38 \pm 0.19$~Ga.

\textbf{15.} $\num{6.75e-4}/\num{8.38e-4} = \qty{0.81}{Ga}$, het is
bijna de helft jonger.

\textbf{16.} $31 + 12 + 7 + 2 + 0 + 2 + 2 = 56$ kraters; $N(\qty{2}{km}) =
56/120000 = \qty{4.67e-4}{km^{-2}}$; $N_1 = \num{4.67e-4}/0.1473 =
\qty{3.17e-3}{km^{-2}}$. Bisectie op $\num{3.22e-14}(e^{6.93T} - 1) +
\num{4.875e-4}\,T = \num{3.17e-3}$ geeft $T = \qty{3.54}{Ga}$ (exponentiële
term \num{1.45e-3}, lineaire term \num{1.73e-3}).

\textbf{17.} $N_1' = \num{3.22e-14}\times6.93\,e^{6.93\times3.54} +
\num{4.875e-4} = \num{1.05e-2}$; $\sigma_T = \num{3.17e-3}/(\sqrt{56}\times
\num{1.05e-2}) = \qty{0.040}{Ga}$.

\textbf{18.} Vlak A: \qty{14}{\percent}; vlak B: \qty{1.1}{\percent}. In
het steile, vroege deel van de chronologie betekent een grote verandering van
de kraterdichtheid een kleine verandering van de ouderdom; op het lineaire
deel is de ouderdom even onzeker als de telling.

\textbf{19.} $120000\times\num{3.17e-3}\times(1 - 0.1473) = 324$; het
draagt $48 + 8 = 56$.

\textbf{20.} $N_1 = 56/(120000\times0.8527) = \qty{5.47e-4}{km^{-2}}$,
$T_r = \num{5.47e-4}/\num{4.875e-4} = \qty{1.12}{Ga}$.

\textbf{21.} $\sigma = 1.12/\sqrt{56} = \qty{0.15}{Ga}$.

\textbf{22.} Een krater van \qty{2}{km} is \qty{400}{m} diep: de lava was
ten minste zo dik. Grotere kraters zijn dieper en hun randen steken hoger
uit; ze zijn deels opgevuld en nog zichtbaar.

\textbf{23.} $1.38 \pm 0.19$ en $1.12 \pm 0.15$~Ga verschillen \qty{0.26}{Ga},
ongeveer één gecombineerde standaarddeviatie ($\sqrt{0.19^2 + 0.15^2} = 0.24$):
de tellingen laten één episode toe, al bewijzen ze die niet.

\textbf{24.} De jonge lavavlakte A heeft een modelouderdom van
\textbf{1.4~miljard~jaar} ($1.38 \pm 0.19$~Ga).
\end{solution}
